% !TeX root = ../navier-stokes-manuscript.tex
% ============================================================
% 2.3 The Pressure-Complete Mixed Jet
% ============================================================

\subsection{The Pressure-Complete Mixed Jet}\label{sub:ThePressureCompleteMixedJet}

Let \(X(a,t)\) be the path of the particle that starts at \(a\), so
\[
  \frac{d}{dt}X(a,t)
  =
  u(X(a,t),t).
\]
Along this path,
\[
  \frac{d}{dt}u(X(a,t),t)
  =
  \partial_tu+(u\cdot\nabla)u
  =:
  D_tu.
\]
The material derivative \(D_tu\) is the particle's acceleration.
The Navier--Stokes equation gives its complete value,
\[
  D_tu
  =
  -\nabla p+\nu\Delta u.
\]
One more carried derivative gives
\begin{align*}
  D_t^2u
  ={}&
  \partial_t^2u
  +2(u\cdot\nabla)\partial_tu
  +(\partial_tu\cdot\nabla)u
  \\
  &+
  (u\cdot\nabla)\bigl((u\cdot\nabla)u\bigr).
\end{align*}
The material jerk is therefore built from Eulerian time derivatives, spatial derivatives, and products of lower responses.
The Eulerian family \(V_n=\partial_t^nu\) gives fixed-coordinate analytic coordinates for these carried responses.

Return to the original system and set
\[
  V_n
  :=
  \partial_t^nu,
  \qquad
  Q_n
  :=
  \partial_t^np.
\]
The undifferentiated row is
\[
  V_1
  +(V_0\cdot\nabla)V_0
  +\nabla Q_0
  =
  \nu\Delta V_0,
  \qquad
  \nabla\cdot V_0=0.
\]
Differentiate once in time.
Leibniz's rule gives
\[
  \partial_t\bigl((V_0\cdot\nabla)V_0\bigr)
  =
  (V_1\cdot\nabla)V_0
  +(V_0\cdot\nabla)V_1,
\]
and hence
\[
  V_2
  +(V_1\cdot\nabla)V_0
  +(V_0\cdot\nabla)V_1
  +\nabla Q_1
  =
  \nu\Delta V_1.
\]
One more differentiation gives
\[
  V_3
  +(V_2\cdot\nabla)V_0
  +2(V_1\cdot\nabla)V_1
  +(V_0\cdot\nabla)V_2
  +\nabla Q_2
  =
  \nu\Delta V_2.
\]
The coefficients \(1,2,1\) appear because two time derivatives can be divided between the transporting and transported factors in those three ways.

At temporal order \(n\), the same product rule gives
\[
  \partial_t^n\bigl((u\cdot\nabla)u\bigr)
  =
  \sum_{a=0}^{n}
  \binom{n}{a}
  (V_a\cdot\nabla)V_{n-a}.
\]
The full temporal recurrence is therefore
\[
  V_{n+1}
  +
  \sum_{a=0}^{n}
  \binom{n}{a}
  (V_a\cdot\nabla)V_{n-a}
  +\nabla Q_n
  =
  \nu\Delta V_n,
  \qquad
  \nabla\cdot V_n=0.
\]
The new time row \(V_{n+1}\) is joined to every lower temporal split of transport, to the pressure response at the same depth, and to two additional spatial derivatives through viscosity.

Pressure is fixed by the same velocity rows.
Taking the divergence of the undifferentiated equation and using \(\nabla\cdot u=0\) gives
\[
  -\Delta p
  =
  \partial_i\partial_j(u_i u_j).
\]
After \(n\) time derivatives,
\[
  -\Delta Q_n
  =
  \partial_i\partial_j
  \sum_{a=0}^{n}
  \binom{n}{a}
  (V_a)_i(V_{n-a})_j.
\]
Rapid decay fixes the additive normalization, equivalently
\[
  Q_n
  =
  R_iR_j
  \sum_{a=0}^{n}
  \binom{n}{a}
  (V_a)_i(V_{n-a})_j.
\]
Thus every temporal pressure row is obtained from the velocity products in the same temporal row.

The recurrence remains open in spatial order, so differentiate it as well.
For a multi-index \(\alpha\in\Nzero^3\), define
\[
  J_{n,\alpha}
  :=
  \partial_t^n\partial_x^\alpha u,
  \qquad
  \Pi_{n,\alpha}
  :=
  \partial_t^n\partial_x^\alpha p.
\]
For \(\beta\leq\alpha\), write
\[
  \binom{\alpha}{\beta}
  :=
  \prod_{k=1}^{3}
  \binom{\alpha_k}{\beta_k}.
\]
The spatial product rule applied to each temporal split gives
\[
  \partial_t^n\partial_x^\alpha
  \bigl((u\cdot\nabla)u\bigr)
  =
  \sum_{a=0}^{n}
  \sum_{\beta\leq\alpha}
  \binom{n}{a}
  \binom{\alpha}{\beta}
  (J_{a,\beta}\cdot\nabla)
  J_{n-a,\alpha-\beta}.
\]
Substitution gives the pressure-complete mixed recurrence
\begin{align*}
  J_{n+1,\alpha}
  &+
  \sum_{a=0}^{n}
  \sum_{\beta\leq\alpha}
  \binom{n}{a}
  \binom{\alpha}{\beta}
  (J_{a,\beta}\cdot\nabla)
  J_{n-a,\alpha-\beta}
  \\
  &+
  \nabla\Pi_{n,\alpha}
  =
  \nu\Delta J_{n,\alpha},
  \qquad
  \nabla\cdot J_{n,\alpha}=0,
\end{align*}
with
\[
  -\Delta\Pi_{n,\alpha}
  =
  \partial_i\partial_j
  \sum_{a=0}^{n}
  \sum_{\beta\leq\alpha}
  \binom{n}{a}
  \binom{\alpha}{\beta}
  (J_{a,\beta})_i
  (J_{n-a,\alpha-\beta})_j.
\]
Every coordinate in this family is a derivative of the same \((u,p)\), and every row retains the transport, viscous, divergence, and pressure relations that produced it.
The mixed recurrence gives the full derivative family.
The critical height now gives a way to see, row by row, how its temporal direction contains an ascending spatial Sobolev column.
